\chapter{Finite quotients and profinite completion}
\label{ch:gr02-profinite}

A finite quotient can distinguish two elements of an infinite group.
The profinite completion puts all of these quotients into one
topological group. The compatible-tuple construction and the product
topology also occur in the Krull-topology comparison in
\emph{Advanced Field Theory} \cite[Results T04--T06]{field}. Here the
indexing objects are the finite-index normal subgroups of an arbitrary
abstract group.

\section{The inverse limit of finite quotients}

\begin{definition}[GR02-C05.01. The profinite completion]
\label{def:gr02-completion}
For an abstract group \(G\), let
\[
  \mathcal N(G)=\{N\trianglelefteq G:[G:N]<\infty\}.
\]
This set includes \(G\), which gives the trivial quotient. For
\(M,N\in\mathcal N(G)\) with \(M\subseteq N\), the transition
homomorphism \(q_{M,N}\colon G/M\to G/N\) sends \(gM\) to \(gN\).
The \emph{profinite completion} of \(G\) is
\[
  \widehat G
  =\left\{(x_N)_{N\in\mathcal N(G)}\in\prod_{N\in\mathcal N(G)}G/N
          :q_{M,N}(x_M)=x_N\text{ if }M\subseteq N\right\}.
\]
We use coordinatewise multiplication and the subspace topology of the
product of finite discrete groups. Write \(p_N\) for coordinate
projection and \(\eta_G(g)=(gN)_N\) for the canonical map.
\end{definition}

The intersection \(M\) of any finite collection of members of
\(\mathcal N(G)\) lies in \(\mathcal N(G)\), since \(G/M\) embeds into
the product of the corresponding finite quotients. In particular a
common refinement lies below each prescribed finite collection under
inclusion. The transition maps compose because each simply passes to
a larger subgroup. Thus these quotients form a cofiltered inverse
system \cite[Definition 1.2.1 and Section 1.2]{profinite-source}.

\begin{proposition}[GR02-C05.02. Topology, density, and kernel]
\label{prop:gr02-completion-topology}
The group \(\widehat G\) is compact, Hausdorff, and totally
disconnected. Each \(p_N\) is a continuous surjection, and the kernels
\(\ker p_N\) form a neighborhood basis at the identity consisting of
clopen normal subgroups. The image of \(\eta_G\) is dense and
\[
     \ker\eta_G=\bigcap_{N\in\mathcal N(G)}N .
\]
\end{proposition}

\begin{proof}
The product of the finite discrete groups is a compact Hausdorff
topological group. Each compatibility condition defines a closed
subgroup, as it is the preimage of the diagonal in
\((G/N)\times(G/N)\) under a continuous group homomorphism. Their
intersection is therefore a compact Hausdorff topological group.
Two distinct tuples have different values in some finite discrete
coordinate, so the clopen fibers of that projection separate them.
This proves total disconnectedness.

Every \(p_N\) is surjective because \(p_N\eta_G\) is the quotient
map. Basic neighborhoods of the identity constrain finitely many
coordinates to be \(1\); the kernel of \(p_M\), for the intersection
\(M\) of those subgroups, lies in the given neighborhood. Each
\(\ker p_N\) is clopen and normal. To approximate any
\(x\in\widehat G\) in a neighborhood prescribing \(x_{N_1},\ldots,x_{N_t}\),
take \(M=N_1\cap\cdots\cap N_t\) and choose \(g\in G\) with
\(gM=x_M\). Compatibility gives \(p_{N_i}\eta_G(g)=x_{N_i}\) for all
\(i\), so every such neighborhood meets \(\eta_G(G)\).
The asserted kernel follows from the formula for \(\eta_G\).
\end{proof}

For a finite \(G\), the coordinate at \(N=\{1\}\) determines all
other coordinates. Thus \(\eta_G\) is an isomorphism onto the
completion, with \(G\) given the discrete topology.

\begin{theorem}[GR02-C05.03. Universal property]
\label{thm:gr02-completion-universal}
Let \(P\) be a profinite group, regarded as a topological inverse
limit of finite discrete groups. Each abstract group homomorphism
\(f\colon G\to P\) has a unique continuous extension
\(\widehat f\colon\widehat G\to P\) satisfying
\(\widehat f\eta_G=f\).
\end{theorem}

\begin{proof}
Write \(P=\varprojlim_i P_i\) with projections \(\pi_i\). For each
\(i\), the kernel \(N_i=\ker(\pi_i f)\) has finite index in \(G\),
since \(G/N_i\) is isomorphic to a subgroup of \(P_i\).
The induced homomorphism \(G/N_i\to P_i\) composed with \(p_{N_i}\)
gives a continuous map \(h_i\colon\widehat G\to P_i\) extending
\(\pi_i f\). For each transition map in the system of the \(P_i\),
the two required coordinate maps agree on \(\eta_G(G)\), by the
compatibility of the coordinates of \(f(g)\). Their equalizer is
closed because the codomain is Hausdorff. Proposition~%
\ref{prop:gr02-completion-topology} therefore implies that they
agree everywhere. The maps \(h_i\) consequently define a
homomorphism to \(P\), continuous in the inverse-limit topology.
On \(\eta_G(G)\) its coordinates are those of \(f\).
Two continuous extensions into the Hausdorff group \(P\) agree on
the same dense subset and are equal.
\end{proof}

Taking \(P=\widehat H\) and composing a group homomorphism
\(G\to H\) with \(\eta_H\) gives its induced continuous homomorphism
on completions. Identities and composition are preserved by
the uniqueness in Theorem~\ref{thm:gr02-completion-universal}.
See \cite[Proposition 3.1.2]{profinite-source} for this comparison.

\begin{proposition}[GR02-C05.04. A cofinal family of finite quotients]
\label{prop:gr02-completion-cofinal}
If \(\mathcal D\subseteq\mathcal N(G)\) has the property that each
\(N\in\mathcal N(G)\) contains some \(D\in\mathcal D\), restriction
induces a topological group isomorphism
\[
       \widehat G \longrightarrow
       \varprojlim_{D\in\mathcal D}G/D .
\]
\end{proposition}

\begin{proof}
For \(D_1,D_2\in\mathcal D\), the condition applied to
\(D_1\cap D_2\) gives a member of \(\mathcal D\) contained in both.
Given a compatible tuple on \(\mathcal D\), define its \(N\)-coordinate
as \(q_{D,N}(x_D)\) from any \(D\subseteq N\) in \(\mathcal D\).
Two choices have a common refinement in \(\mathcal D\), which
shows that this coordinate is independent of the choice.
If \(M\subseteq N\), choosing \(D\subseteq M\) computes both
coordinates and proves their compatibility. Each coordinate of
this extension is a continuous group homomorphism since it is
given by one projection followed by a finite quotient map.
The extension and restriction recover the original coordinates
in both directions.
\end{proof}

\begin{example}[GR02-C05.05. The integers and finite fields]
\label{ex:gr02-completion-integers}
The finite-index subgroups of \(\mathbb Z\) are \(n\mathbb Z\) for
positive integers \(n\). The coordinates of the completion are
therefore the compatible residue classes in
\(\varprojlim_n\mathbb Z/n\mathbb Z\), with transition from \(m\)
to \(n\) when \(n\mid m\). One obtains the same limit using
the quotients by \(j!\mathbb Z\), since every positive integer
divides some factorial. Proposition~\ref{prop:gr02-completion-cofinal}
identifies the two topologies as well as the groups.

The absolute Galois group of \(\mathbb F_q\) is topologically
isomorphic to this completion, with \(x\mapsto x^q\) corresponding
to \(\eta_{\mathbb Z}(1)\). This is the residue comparison proved
in \emph{Advanced Field Theory} \cite[Result E06]{field}.
\end{example}

\section{Residual finiteness and three examples}

\begin{proposition}[GR02-C05.06. Residual finiteness and injectivity]
\label{prop:gr02-residual-finiteness}
An abstract group \(G\) is \emph{residually finite} if for each
\(g\ne 1\) some homomorphism to a finite group has nonidentity
value on \(g\). This is equivalent to \(\eta_G\) being injective.
In a residually finite group every finite subset maps injectively
to a single finite quotient.
\end{proposition}

\begin{proof}
A homomorphism to a finite group factors through the finite quotient
by its kernel. Conversely each quotient by \(N\in\mathcal N(G)\)
provides such a homomorphism. Thus residual finiteness is equivalent
to the intersection in Proposition~\ref{prop:gr02-completion-topology}
being trivial. Given a finite subset, choose a finite quotient
that detects \(x^{-1}y\) for each pair of distinct elements \(x,y\)
in it. The intersection of the finitely many kernels gives one
finite quotient in which all these pairs remain distinct.
For a subset of size at most one the trivial quotient suffices.
\end{proof}

\begin{proposition}[GR02-C05.07. Finitely generated abelian groups]
\label{prop:gr02-abelian-residual}
Every finitely generated abelian group is residually finite.
\end{proposition}

\begin{proof}
By the structure theorem \cite[K07.07.fg\_abelian\_extension]{d07},
such a group is isomorphic to \(\mathbb Z^r\oplus A\) with \(A\)
finite abelian. For a nonzero element with nonzero \(A\)-component,
projection to \(A\) detects it. Otherwise at least one integer
coordinate \(a\) is nonzero. Project to that coordinate and reduce
modulo \(|a|+1\), so that the residue of \(a\) is nonzero.
\end{proof}

\begin{theorem}[GR02-C05.08. Free groups]
\label{thm:gr02-free-residual}
A free group on any set of generators is residually finite.
\end{theorem}

\begin{proof}
By the reduced-word description of the free group
\cite[Result A01]{geomgroup}, a nonidentity element has a nonempty
reduced word \(w=s_{\ell}^{\varepsilon_{\ell}}\cdots
s_1^{\varepsilon_1}\), where each \(\varepsilon_i\) is \(1\) or
\(-1\) and each \(s_i\) is a generator. Use permutations acting
on the left, so the rightmost letter acts first. On
\(X=\{0,\ldots,\ell\}\), prescribe part of the permutation
\(\sigma_s\) for each generator \(s\) that occurs in \(w\):
\[
   \sigma_s(i-1)=i\quad\text{if }s_i=s,\ \varepsilon_i=1,\qquad
   \sigma_s(i)=i-1\quad\text{if }s_i=s,\ \varepsilon_i=-1.
\]
Two different steps for \(s\) could prescribe the same domain
only at a shared vertex of consecutive steps, with signs
\(-1,1\) in step order. Two different steps could prescribe the
same range only at a shared vertex with signs \(1,-1\).
Either would put adjacent inverse letters in \(w\).
Thus each prescription is a bijection between two subsets of \(X\).
Their complements have equal finite cardinalities, so choose a
bijection between the complements to extend it to a permutation.
For generators absent from the word, take the identity.

By the free-group universal property these assignments define a
homomorphism to the finite permutation group of \(X\).
For each \(i\), the letter \(s_i^{\varepsilon_i}\) moves
\(i-1\) to \(i\) by construction. The image of \(w\) consequently
moves \(0\) to \(\ell\ne 0\). This finite quotient (taking the
image of the homomorphism if necessary) detects \(w\).
The construction uses only the letters present in the word.
\end{proof}

Completing these partial permutations amounts to completing the
labelled path to a finite covering graph, as in
\cite[Theorem 3.2.1, first proof]{profinite-source}.

\begin{example}[GR02-C05.09. The additive rationals]
\label{ex:gr02-rationals-completion}
If \(f\colon(\mathbb Q,+)\to H\) has finite target of order \(n\),
then for every \(x\in\mathbb Q\)
\[
   f(x)=f(x/n)^n=1.
\]
The last equality follows from Lagrange's theorem in \(H\).
All the finite quotients are therefore trivial and so is
\(\widehat{\mathbb Q}\). Its canonical map still has dense image,
as in Proposition~\ref{prop:gr02-completion-topology}.
\end{example}
